\section{证据边界与术语消化}

本讲跨越架构、系统与产品结果，容易产生两个误解：一是用单张样例替代全面评估，二是用 FLOPs 替代真实延迟。Q\&A 对 GAN、自动指标、视频成本、AdaLN 与 MMDiT 的追问，提供了必要的工程边界。

\subsection{Q\&A 带来的四条校正}

第一，GAN 没有被彻底淘汰：当 one-shot、超低延迟比最高质量更重要时仍有价值。第二，图像/视频自动指标很弱，必须结合人类偏好与任务指标。第三，视频不是独立帧集合。第四，MMDiT 的直觉不是理论最优证明。

\teachervoice{00:57:31--01:01:38，在线问答讨论 autoregressive 模型、GAN、自动指标与视频挑战。老师没有给“扩散统一胜出”的答案，而是按 latency、quality 与 temporal consistency 区分场景。}

\teachervoice{01:01:45--01:03:00，老师再次解释 AdaLN：LayerNorm 负责稳定特征，adaptive 部分用外部条件改变 scale/shift，从而让同一 block 在不同 timestep 或条件下执行不同计算。}

\begin{warningbox}{Qualitative cherry-picking}
生成模型可以从大量样本中挑出少数最佳图。可靠报告应固定 prompt set、公开 seed、展示失败样本，并同时测 prompt fidelity、diversity、human preference、latency 与 memory。
\end{warningbox}

\subsection{术语集中消化}

下面把课程中的高密度术语按“解决的问题”重新组织，避免把模型名当成机制。

\begin{center}
\small
\begin{tabularx}{\textwidth}{p{0.18\textwidth} X X}
\toprule
\textbf{术语} & \textbf{解决的问题} & \textbf{核心机制与课程关系} \\
\midrule
Scheduler & 如何从一个噪声状态走到下一个 & 定义 timestep、噪声日程和数值更新；不等于 denoiser。 \\
Denoiser & 每一步应预测什么方向 & U-Net 或 DiT，输出 noise、velocity 或 data parameterization。 \\
Latent diffusion & 如何降低像素空间成本 & 先用 VAE 压缩，再在低分辨率 latent 上迭代。 \\
U-Net & 如何利用多尺度图像偏置 & Down/up blocks 与 skip connections 保留局部和全局信息。 \\
DiT & 如何用标准 Transformer 做去噪 & Patchify latent，用 timestep/condition 调制 encoder blocks。 \\
AdaLN-Zero & 如何低成本注入条件并稳定深网络 & 条件生成 scale/shift/gates，零初始化使 block 初始近似 identity。 \\
Cross-attention & 图像 token 如何读取文本 & 图像作 query，文本作 key/value；通常只更新图像流。 \\
Linear attention & 如何避免 image self-attention 的 $N^2$ & 重排核化 attention 的乘法次序，牺牲部分精确交互。 \\
MMDiT & 两种模态如何保留专属空间又共同演化 & 分开 AdaLN/QKV，拼接后做 joint attention。 \\
QK-Norm & 大模型 attention logits 如何稳定 & 规范化 query/key，控制点积尺度。 \\
GQA & 如何降低 key/value 成本 & 多个 query heads 共享较少 K/V heads。 \\
RoPE & 如何编码相对空间与时间位置 & 旋转 query/key，可扩展到视频三轴。 \\
Structural control & 如何遵循 edge/pose/depth & 编码空间对齐条件并注入 base model。 \\
Parameter sharing & 如何减少权重与部分计算 & 跨层或跨模态共享 AdaLN、QKVO、MLP，形成质量--效率折中。 \\
\bottomrule
\end{tabularx}
\end{center}

\subsection{本章小结}

生成架构不能脱离目标与 workload 比较。GAN、U-Net、DiT、linear attention 与 MMDiT 各有适用区间；自动指标、参数量与 FLOPs 都只是局部代理。可信结论需要控制变量、端到端测量和失败样本。

